\section{Experimental methodology and demonstrator testbed}\label{sec:methodology}

This section describes the planned demonstrator and the integrated experimental and numerical methodology through which the Self-Aware Steel Structure concept will be instantiated and assessed. No experiments have yet been performed; what follows is the proposed protocol. The strategy deliberately couples a physical testbed, a continuously updated finite-element representation and the machine-learning layer described in the preceding sections, so that the structure is not merely instrumented but is enabled to interpret and report its own mechanical state.

\subsection{Physical demonstrator}\label{subsec:demonstrator}

The demonstrator will be a three-dimensional braced frame fabricated from square hollow structural steel sections (SHS). This configuration is chosen because it is representative of skeletal steel systems widely used in industrial buildings, warehouses, elevated platforms and modular construction, where bar- and beam-type idealizations are appropriate \citep{cook2002fea}. Beam-to-column connections will be bolted and idealized as pinned, while global lateral stability will be provided exclusively by diagonal bracing. The arrangement isolates two well-defined and physically meaningful damage scenarios, loss of bracing axial stiffness and loosening of bolted connections, both of which manifest as changes in material or geometric properties, boundary conditions or system connectivity, and therefore satisfy the formal definition of damage as a change of state adversely affecting performance \citep{farrar2013shm}. A loosened bolted connection, in particular, alters connectivity and typically increases energy dissipation, producing a measurable rise in vibrational damping that the monitoring system is designed to capture. The complete instrumented demonstrator and its coupling to the digital-twin loop are shown in Fig.~\ref{fig:testbed}.

\subsection{Instrumentation and sensor-placement strategy}\label{subsec:instrumentation}

A distributed, heterogeneous sensor network will be deployed. The planned channels comprise electrical resistance strain gauges (local strains and, via the constitutive relation, stresses); triaxial accelerometers (dynamic response for modal identification); temperature sensors (thermal compensation and operational or environmental normalization); displacement transducers (static and quasi-static deflections); and inclinometers (global stability and drift). Sensor positions will not be chosen ad hoc but will be selected from a preliminary numerical analysis of the FEM. Candidate locations will be ranked by their sensitivity to the targeted modal quantities and damage scenarios, so that accelerometers are placed where the lower mode shapes exhibit large modal amplitude and strain gauges where the damage-sensitive members carry significant force. This model-informed placement reflects the principle that sensing-system requirements are dictated by the length and time scales of the damage of interest, and that distributed arrays plus heterogeneous data fusion improve detection fidelity \citep{farrar2013shm}. It is also consistent with established computational sensor-placement strategies and optimization algorithms developed for SHM \citep{tan2020sensorplacement,hassani2023sensorplacement}. The full channel list, the measured quantity and the placement rationale for each sensor type are summarized in Table~\ref{tab:sensors}. All signals will be continuously acquired, synchronized and time-stamped by a central data-acquisition (DAQ) system, supplying the streams of strain, acceleration, displacement, temperature and inclination from which the higher-level features, natural frequencies, mode shapes, damping ratios and stiffness indicators, will be derived.

\subsection{Excitation system and frequency sweeps}\label{subsec:excitation}

Controlled dynamic loading will be applied by a rotating-unbalance excitation system: a variable-speed electric motor driving an eccentric mass rigidly attached to the frame. This is the textbook mechanism for applying harmonic excitation to full-scale structures via rotating machinery \citep{cook2002fea}, with a force amplitude proportional to the square of the excitation frequency, $p(t)=m_e\,e\,\omega^2\sin\omega t$ \citep{chopra2017dynamics}. Because the force scales with $\omega^2$, the device is well suited to dynamic rather than static testing. The motor speed will be swept across a range that brackets the lower natural frequencies of the frame, allowing the steady-state frequency-response curve to be traced; its peaks locate the resonant frequencies and their sharpness encodes damping through the half-power bandwidth \citep{chopra2017dynamics}. By exciting at and near resonance, the test campaign will deliberately amplify the response in order to reveal modal shifts, resonance migration and changes in resonant amplitude induced by introduced stiffness changes or damage. The physical basis is the direct coupling $\omega_n=\sqrt{k/m}$, whereby any stiffness loss is observable as a downward shift in natural frequency and an increase in measured damping \citep{chopra2017dynamics}.

\subsection{Finite-element model and parameter identification}\label{subsec:femodel}

A finite-element model will reproduce the static and dynamic behavior of the demonstrator. SHS members will be discretized with beam/frame elements assembled into the global stiffness matrix $[\mathbf{K}]$; mass will be represented by a consistent (or, where appropriate, lumped) mass matrix $[\mathbf{M}]$, and energy dissipation by Rayleigh proportional damping $[\mathbf{C}]=a[\mathbf{M}]+b[\mathbf{K}]$ \citep{cook2002fea}. The model will incorporate section geometry, measured material properties, idealized boundary conditions and connection flexibility, with the pinned or rigid idealization of bolted joints treated explicitly as a modeling decision to be reconciled with measurement. Four analysis types will establish the baseline: static analysis $[\mathbf{K}]\{\mathbf{D}\}=\{\mathbf{R}\}$; modal analysis via the undamped eigenproblem $([\mathbf{K}]-\omega_n^2[\mathbf{M}])\{\boldsymbol{\phi}\}=\{\mathbf{0}\}$, yielding the baseline natural frequencies and mode shapes; harmonic (frequency-response) analysis to predict steady-state amplitude versus excitation frequency under the rotating-unbalance load; and transient analysis for non-stationary events \citep{cook2002fea,chopra2017dynamics}. Recognizing that finite-element analysis is simulation, not reality \citep{cook2002fea}, the model will be calibrated against the undamaged structure: experimentally identified modal parameters will be used to update stiffness, connection-flexibility and damping parameters so that the numerical baseline faithfully represents the as-built healthy state. Comparable modal identification and system-identification campaigns on full-scale steel members support the feasibility of this calibration step \citep{yu2024steelbridge}.

\subsection{Digital-twin integration and deviation indicators}\label{subsec:dtintegration}

The validated FEM will serve as the kernel of a digital twin continuously fed by the live sensor streams \citep{chiachio2022structural}. At each monitoring step, measured responses (strains, frequencies, mode shapes, damping, displacements) will be compared in real time against the twin's predictions for the current operating and environmental condition, with temperature channels used to normalize benign variability so that environmental effects are not mistaken for damage \citep{farrar2013shm}. Deviations will be quantified through statistical indicators and performance parameters: residuals between measured and predicted features, control-chart-style thresholds, and distance metrics on the modal and feature space, providing the quantitative bridge between physical state and digital representation \citep{sohn2003statistical}. Persistent, statistically significant deviations from the healthy baseline will then be passed to the machine-learning layer for novelty detection, classification and severity assessment, closing the perceive--process--interpret--report loop that defines Structural Cognition \citep{zou2026deep}.

\subsection{Six-step methodology}\label{subsec:sixstep}

The study will proceed through the following sequential and partly iterative steps:

\begin{enumerate}
    \item \textbf{Develop the finite-element model} of the SHS braced frame, including geometry, materials, connection flexibility and boundary conditions, and run static, modal, harmonic and transient analyses to establish baseline behavior.
    \item \textbf{Build the instrumented physical structure}, fabricating and erecting the bolted, braced 3D frame as the experimental demonstrator.
    \item \textbf{Install and calibrate the distributed sensor network} (strain gauges, accelerometers, temperature sensors, displacement transducers, inclinometers) at the numerically selected positions, and integrate the central DAQ.
    \item \textbf{Run static and dynamic tests}, including rotating-unbalance frequency sweeps near the natural frequencies, on the healthy structure and under controlled, reversible damage scenarios (bracing stiffness reduction, bolted-connection loosening).
    \item \textbf{Implement the digital twin and integrate real-time data}, performing continuous FEM-versus-measurement comparison with statistical deviation indicators and environmental normalization.
    \item \textbf{Train, validate and apply the AI layer} (supervised and unsupervised models) for autonomous structural-integrity assessment, mapping deviations to the existence, location, type and severity of damage.
\end{enumerate}

To make this plan auditable rather than merely descriptive, Table~\ref{tab:params} consolidates the parameters that fully define the campaign: geometry, material, connections, sensor count and placement, sampling and filtering, excitation, quantified damage levels, the model-updating procedure, the alarm thresholds, and the train/validation/test partition. These values are deliberately presented as a specification to be finalised with the as-built structure; the present paper fixes the protocol and the quantities to be reported, not measured results.
